\subsection{The Constructor Comment}\label{sec:ConstructorComment}



\}


Basically, a constructor is a special kind of a method, so the same rules are valid for the documentation comment for a constructor than for the documentation comment for a method, as given in chapter \tqref{sec:MethodComment}.

For an (empty) default constructor, the constructor comment may be as simple as this, no matter if it is \lstinline|public|, \lstinline|private|, or \lstinline|protected|:
\begin{lstlisting}
/**
 *  Creates a new instance of {@code MyClass}.
 */
public MyClass() { /* Does nothing */ }
\end{lstlisting}
or
\begin{lstlisting}
/**
 *  Default constructor for class {@code MyClass}.
 */
private MyClass() { /* Does nothing */ }
\end{lstlisting}
with the first alternative being the preferred one, because this comment can be used for every constructor, not only for the default ones.

Of course, this comment do not say very much, but the constructor of a class is also not doing that much, and what it does is very obvious. But if the constructor has side effects, these should be described properly.

A class that does have only \lstinline|static| methods\footnote{Such a class is called a “Utility Class”; refer to \tqfullvref{sec:UtilityClasses} for more details.} should have a \lstinline|private| constructor like this:\footnote{The class for the \lstinline|Error| used in the \lstinline|throw| is described in \autocite{TQUADRAT_ORG_FOUNDATION_PRIVATECONSTRUCTORFORSTATICCLASSCALLEDERROR}.}
\begin{lstlisting}
/**
 *  No instance is allowed for class {@code MyUtilityClass}.
 */
private MyUtilityClass() 
{ 
    throw new PrivateConstructorForStaticClassCalledError( MyUtilityClass.class ); 
}   // MyUtilityClass()
\end{lstlisting}

If a constructor takes parameters, there has to be a \nameref{sec:TagParam} tag for each of them, exactly like for a method.

Although constructors should not throw (checked) exceptions, sometimes it could not be avoided without overcomplicating the API of a class. In such case, all exceptions has to be listed with the \nameref{sec:TagThrows} tag and a description of the conditions for the particular exception – as far as it is possible or make sense. Refer also to chapter \tqfullvref{sec:GeneralExceptionHandling} for some more details on exception handling.

Same as for a method, the documentation comment for a constructor has to describe all the side-effects caused by it; the modification of an argument is such a side-effect – refer to chapter \tqfullvref{sec:MutableTypesForArguments} for more details. But in general, a constructor should not have any (visible) side-effects at all.


%==============================================================================
% End of file!!
%==============================================================================

